\clearpage
\begin{proofbox}
    \textbf{\textcolor{CProof}{方法一：门槛最低——圆的切线与勾股定理}}

\begin{center}
\begin{tikzpicture}[x=2.5cm,y=2.5cm,font=\small,
  every node/.style={inner sep=3pt},line cap=round]

\coordinate (O) at (0,0); \coordinate (T) at (0,1);
\coordinate (Q) at (1.333333,1);
\draw[->,gray!65] (-1.3,0)--(2.1,0) node[right] {$u$};
\draw[->,gray!65] (0,-1.3)--(0,1.55) node[above] {$v$};
\draw[CStatement,very thick] (O) circle (1);
\fill[CProof!12] (O)--(T)--(Q)--cycle;
\draw[CProof,thick] (O)--(T) node[midway,left] {$1$};
\draw[orange!85!black,very thick] (-.35,1)--(1.7,1);
\draw[CProof,thick] (O)--(Q);
\draw (0,.88)--(.12,.88)--(.12,1);
\fill (O) circle (.8pt) node[below left] {$O$};
\fill (T) circle (.8pt) node[above left] {$T(u_1,v_1)$};
\fill (Q) circle (.8pt) node[above right] {$Q(u_0,v_0)$};

\end{tikzpicture}

{\small 圆中 $OT\perp TQ$，所以 $OQ^2=1+TQ^2$；直角只标在圆中。}
\end{center}


    仍设 $a,b>0$，外点 $P(x_0,y_0)$ 的两个不同切点为 $A,B$。

    \textbf{核心：} 把椭圆缩放成单位圆，把问题转成初中熟悉的``半径垂直切线''。
    \begin{description}[style=nextline, leftmargin=2.8em, labelsep=0.5em, itemsep=0.4em]
        \item[\textbf{第1步（缩放成圆）：}] 令 $u=\frac xa,\ v=\frac yb$，椭圆变为 $u^2+v^2=1$，外点
        $P$ 变为 $Q(u_0,v_0)=(\frac{x_0}{a},\frac{y_0}{b})$。
        横纵坐标按固定比例分别缩放，直线仍是直线；并且对应的相交点一一对应，所以原来的切线映成圆的切线。

        \item[\textbf{第2步（使用勾股定理）：}] 设圆上一切点为 $T(u_1,v_1)$。
        圆心 $O$ 到切点的半径 $OT\perp QT$，且 $OT=1$。因此
        \[
            u_0^2+v_0^2=1+(u_0-u_1)^2+(v_0-v_1)^2.
        \]
        展开并利用 $u_1^2+v_1^2=1$，化简得
        \[
            \boxed{u_0u_1+v_0v_1=1}.
        \]
        \textbf{理由：} 这是把直角三角形中的三边分别用两点间距离表示。

        \item[\textbf{第3步（换回椭圆坐标）：}] 把
        $u_0=\frac{x_0}{a},\ v_0=\frac{y_0}{b},\ u_1=\frac{x_1}{a},\ v_1=\frac{y_1}{b}$
        代入，得
        \[
            \frac{x_0x_1}{a^2}+\frac{y_0y_1}{b^2}=1.
        \]
        两个切点 $A,B$ 都满足这同一个一次方程，所以它就是经过 $A,B$ 的直线方程。\textbf{证毕。}
    \end{description}

\end{proofbox}
\clearpage
\begin{proofbox}
    \textbf{\textcolor{CProof}{方法二：证明最快——直接利用已知切线公式}}

\begin{center}
\begin{tikzpicture}[x=1.28cm,y=1.28cm,font=\small,
  every node/.style={inner sep=3pt},line cap=round]

\coordinate (O) at (0,0); \coordinate (A) at (0,2);
\coordinate (B) at (2.88,-.56); \coordinate (P) at (4,2);
\draw[->,gray!65] (-3.4,0)--(4.6,0) node[right] {$x$};
\draw[->,gray!65] (0,-2.4)--(0,2.7) node[above] {$y$};
\draw[CStatement,very thick] (O) ellipse [x radius=3,y radius=2];
\draw[orange!85!black,thick] (-.65,2)--(4.35,2);
\draw[orange!85!black,thick] (2.55,-1.3143)--(4,2);
\draw[CTrap,very thick] (A)--(B);
\fill (O) circle (1.4pt) node[below left] {$O$};
\fill (A) circle (1.7pt) node[above left] {$A$};
\fill (B) circle (1.7pt) node[below right] {$T=B$};
\fill (P) circle (1.7pt) node[above right] {$P(x_0,y_0)$};

\node[anchor=west,text=CProof] at (-3.15,-2.85) {$T(x_1,y_1)$ 的切线：$\dfrac{x_1x}{a^2}+\dfrac{y_1y}{b^2}=1$};

\end{tikzpicture}

{\small 任一切点 $T$ 的切线经过 $P$；两切点都满足同一乘积式。}
\end{center}


    \textbf{适用前提：} 已掌握椭圆上一点 $(x_1,y_1)$ 处的切线方程
    $\frac{x_1x}{a^2}+\frac{y_1y}{b^2}=1$。
    \begin{description}[style=nextline, leftmargin=2.8em, labelsep=0.5em, itemsep=0.35em]
        \item[\textbf{第1步：}] 任取一个切点 $T(x_1,y_1)$，其切线经过 $P$，故
        \[
            \frac{x_1x_0}{a^2}+\frac{y_1y_0}{b^2}=1.
        \]
        \item[\textbf{第2步：}] 交换乘法次序：$\frac{x_0x_1}{a^2}+\frac{y_0y_1}{b^2}=1$。
        所以任一切点均在直线 $\frac{x_0x}{a^2}+\frac{y_0y}{b^2}=1$ 上。
        \item[\textbf{第3步：}] 两个不同切点确定直线 $AB$，故所求即
        $\boxed{\frac{x_0x}{a^2}+\frac{y_0y}{b^2}=1}$。\textbf{证毕。}
    \end{description}
    \textbf{考场速写：} ``切点 $T$ 的切线过 $P$ $\Rightarrow$ $T$ 满足乘积式 $\Rightarrow$ 两切点共线。''

\end{proofbox}
\clearpage
\begin{proofbox}
    \textbf{\textcolor{CProof}{方法三：常规通用——联立直线与椭圆，令交点只有一个}}

\begin{center}
\begin{tikzpicture}[x=1.28cm,y=1.28cm,font=\small,
  every node/.style={inner sep=3pt},line cap=round]

\draw[->,gray!65] (-3.4,0)--(4.8,0) node[right] {$x$};
\draw[->,gray!65] (0,-2.35)--(0,3.25) node[above] {$y$};
\draw[CStatement,very thick] (0,0) ellipse [x radius=3,y radius=2];
\draw[orange!85!black,very thick,->] (-2.8,2)--(4.6,2);
\fill (0,2) circle (1.7pt) node[above left] {$T(x_1,y_1),\ t=0$};
\fill (4,2) circle (1.7pt) node[above] {$P,\ t=1$};
\fill (2,2) circle (1.7pt) node[above,yshift=1pt] {$t=\frac12$};
\draw[->,CProof,thick] (.6,2.65)--(3.1,2.65) node[midway,above] {$t\text{ 增大}$};
\node at (0,-.3) {$O$};

\end{tikzpicture}

{\small 直线 $PT$：$T$ 对应 $t=0$，$P$ 对应 $t=1$；相切时只有 $T$ 在椭圆上。}
\end{center}


    \textbf{特点：} 不预先使用椭圆切线公式，也不需要分讨论直线斜率是否存在。
    \begin{description}[style=nextline, leftmargin=2.8em, labelsep=0.5em, itemsep=0.45em]
        \item[\textbf{第1步（设切点和直线）：}] 设 $T(x_1,y_1)$ 是其中一个切点，故
        \[
            \frac{x_1^2}{a^2}+\frac{y_1^2}{b^2}=1.
        \]
        因为直线 $PT$ 经过两点 $T,P$，其上任一点可写作
        \[
            x=x_1+t(x_0-x_1),\quad y=y_1+t(y_0-y_1),
        \]
        其中 $t$ 为实数；$t=0$ 对应 $T$，$t=1$ 对应 $P$。

        \item[\textbf{第2步（代入椭圆并展开）：}] 将上式代入椭圆方程：
        \[
        \begin{aligned}
           &\frac{[x_1+t(x_0-x_1)]^2}{a^2}
             +\frac{[y_1+t(y_0-y_1)]^2}{b^2}=1,\\
           &\Rightarrow\quad
           t^2\!\left(\frac{(x_0-x_1)^2}{a^2}+\frac{(y_0-y_1)^2}{b^2}\right)
           +2t\!\left(\frac{x_1(x_0-x_1)}{a^2}+\frac{y_1(y_0-y_1)}{b^2}\right)=0.
        \end{aligned}
        \]
        常数项消失，\textbf{因为 $T$ 本来就在椭圆上}。

        \item[\textbf{第3步（相切意味着只有一个交点）：}] 上式已经有解 $t=0$。
        因为 $PT$ 是切线，不能有另一个不同的交点，所以二次方程的两个根都必须是 $0$。
        又因 $P\ne T$，$t^2$ 的系数大于 $0$，故一次项的系数为 $0$：
        \[
            \frac{x_1(x_0-x_1)}{a^2}+\frac{y_1(y_0-y_1)}{b^2}=0.
        \]
        移项并用 $\frac{x_1^2}{a^2}+\frac{y_1^2}{b^2}=1$，得到
        \[
            \boxed{\frac{x_0x_1}{a^2}+\frac{y_0y_1}{b^2}=1}.
        \]

        \item[\textbf{第4步（确定切点弦）：}] 这个推导对两个切点 $A,B$ 分别成立，因此两点都在
        $\frac{x_0x}{a^2}+\frac{y_0y}{b^2}=1$ 上；该直线就是 $AB$。\textbf{证毕。}
    \end{description}

    \medskip
    \textbf{\textcolor{CProof}{三法对照}}

    \textbf{方法一} 几何最直观，突出``勾股定理 $\rightarrow$ 坐标乘积''；
    \textbf{方法二} 最省时，建立在已会椭圆切线公式的基础上；
    \textbf{方法三} 代数可复查，能够独立说明椭圆切线公式为何成立，竖直切线也不漏掉。
\end{proofbox}
